\section{Unitary lift of refinement and transport of readers}
\label{lib05:levantamiento}

The residue bijection of \cref{lib03:teo-residuo} and the
bilateral balance of \cref{lib04:teo-balance} admit an
explicit composition. An admissible arithmetic triple, containing
a pair and its carry, determines a refined pair. The
linearization of this bijection on orthonormal bases provides
a unitary transport; two such transports can act
as the two realizations of the bilateral balance.

The realization is established downstream of the objects generated by
APP--TRIT--TPK. Its combinatorial domain comprises the pairs and
carries satisfying the conditions stated below. The norm
preserved by its linearization is the sum of squared amplitudes
over those labels; the arithmetic conservation law for
\(L_c(x,y)=(Bx-c,By+c)\) remains the sum of the two
coordinates after rescaling. The composition relates the two
conservation laws through a concrete map.

\subsection{Bijection of admissible arithmetic states}
\label{lib05:biyeccion}

Let \(B\geq2\) and \(M\geq1\) be integers. Define
\begin{align}
 \mathcal D_{B,M}
 &=\bigl\{(x,y,c)\in\mathbb Z_{\geq0}^2
                 \times\{0,\ldots,B-1\}:x+y=M,\ c\leq Bx\bigr\},
 \label{lib05:dominio}\\
 \mathcal E_{BM}
 &=\bigl\{(X,Y)\in\mathbb Z_{\geq0}^2:X+Y=BM\bigr\}.
 \label{lib05:codominio}
\end{align}
The condition \(c\leq Bx\) preserves the nonnegativity of the
first output coordinate. For \(x=0\), it permits only
\(c=0\); for \(x\geq1\), it permits the entire canonical alphabet.

\begin{theorem}[Exact correspondence of labels]
\label{lib05:teo-biyeccion}
The map
\begin{equation}
 \ell_{B,M}:\mathcal D_{B,M}\longrightarrow\mathcal E_{BM},
 \qquad (x,y,c)\longmapsto(Bx-c,By+c)
 \label{lib05:mapa-etiquetas}
\end{equation}
is bijective. Both sets have cardinality \(BM+1\). Its inverse
is given by
\begin{equation}
 c=Y\bmod B,\qquad y=\frac{Y-c}{B},\qquad x=\frac{X+c}{B}.
 \label{lib05:inversa-etiquetas}
\end{equation}
\end{theorem}

\begin{proof}
The domain condition and the identity
\(Bx-c+By+c=BM\) ensure that the image belongs to
\(\mathcal E_{BM}\). Given a pair in this set, the residue
\(c\) is the unique element of the alphabet satisfying
\(Y\equiv c\pmod B\). Since \(X+Y\) is divisible by \(B\),
so is \(X+c\). Euclidean division of \(Y\) gives
\(y\geq0\), and \(X+c\geq0\) gives \(x\geq0\). If \(c>0\),
\(x=0\) would require \(X=-c<0\), so admissibility
is preserved. The preimage coordinates sum to \(M\), and
substitution verifies the original image. Uniqueness of the residue and
quotients proves the bijection.

There are \(M+1\) labels with \(c=0\) and \(M\) with each positive
carry; hence the cardinality is
\((M+1)+(B-1)M=BM+1\), equal to that of the output.
\end{proof}

The proof makes the boundary endpoint of the domain explicit, a condition
necessary for equality of cardinalities. The new carry label
is a typed part of the input to refinement.

\subsection{Unitarity on label amplitudes}
\label{lib05:unitariedad}

Take the real or complex Hilbert spaces with counting measure
\begin{equation}
 \mathscr H_D=\ell^2(\mathcal D_{B,M}),\qquad
 \mathscr H_E=\ell^2(\mathcal E_{BM}),
 \label{lib05:hilbert}
\end{equation}
and their labelled orthonormal bases. Define linearly
\begin{equation}
 J_{B,M}|x,y,c\rangle=|Bx-c,By+c\rangle.
 \label{lib05:j}
\end{equation}
We write \(J\) when the parameters are fixed.

\begin{proposition}[Unitary lift]
\label{lib05:prop-unitaria}
The operator \(J:\mathscr H_D\to\mathscr H_E\) is unitary:
\(J^*J=I_D\), \(JJ^*=I_E\). The adjoint recovers each label
through \eqref{lib05:inversa-etiquetas}.
\end{proposition}

\begin{proof}
The result in \cref{lib05:teo-biyeccion} ensures that \(J\) maps
two complete orthonormal bases bijectively onto one another. For
\(\psi=\sum_{d\in\mathcal D}\psi_d|d\rangle\),
\[
 \|J\psi\|^2
 =\sum_{e\in\mathcal E}|\psi_{\ell_{B,M}^{-1}(e)}|^2
 =\sum_{d\in\mathcal D}|\psi_d|^2=\|\psi\|^2.
\]
Surjectivity proves the two unitary identities, and the
inverse bijection describes the adjoint on each basis vector.
\end{proof}

The label \(|73,27,1\rangle\) has norm one; the integers
\(73,27,1\) describe its arithmetic values, not its amplitudes.
The construction represents reversible processing of labels
or amplitudes. A concrete physical realization also retains
the protocol that prepares and transports those amplitudes.

\subsection{Diagonal readers of proportion and carry}
\label{lib05:diagonales}

For a real-valued function \(h\) on a finite set, let
\(M_h\) denote its diagonal operator: each basis vector is multiplied
by the value of \(h\) at its label. The bijection satisfies
\begin{equation}
 J^*M_hJ=M_{h\circ\ell_{B,M}}.
 \label{lib05:lector-diagonal}
\end{equation}
This follows by acting on every basis vector.

Define on the output
\begin{equation}
 f_E(X,Y)=\frac{X}{BM},\quad
 g_E(X,Y)=\frac{Y}{BM},\quad
 \gamma_E(X,Y)=Y\bmod B,
 \label{lib05:lectores-salida}
\end{equation}
and on the input \(f_D=x/M\), \(g_D=y/M\), \(\gamma_D=c\).
Identity \eqref{lib05:lector-diagonal} yields
\begin{align}
 J^*M_{f_E}J&=M_{\,x/M-c/(BM)},\nonumber\\
 J^*M_{g_E}J&=M_{\,y/M+c/(BM)},\nonumber\\
 J^*M_{\gamma_E}J&=M_c.
 \label{lib05:transporte-fracciones}
\end{align}
The first two readers transport the transfer of a unit.
Since \(f_E+g_E=1\), their sum is \(I_E\) and its compression is
\(I_D\). The third reader preserves the carry exactly as an
integer observable, with its residue representation in the output.

\subsection{Complete paths at arbitrary depth}
\label{lib05:rutas}

Let \(\mathcal P_n(B,M)\) be the set of paths of length \(n\)
starting from all nonnegative pairs with sum \(M\) and, at
each stage, satisfying the conditions of
\eqref{lib05:dominio} for the corresponding sum. A path retains
its initial pair and its \(n\) admissible carries. Let
\(\Phi_n\) be the map assigning to each path its final pair.

\begin{theorem}[Recovery of paths at every finite depth]
\label{lib05:teo-rutas}
For every \(n\geq0\),
\begin{equation}
 \Phi_n:\mathcal P_n(B,M)\longrightarrow\mathcal E_{B^nM}
 \quad\hbox{is bijective},\qquad
 |\mathcal P_n(B,M)|=B^nM+1.
 \label{lib05:biyeccion-rutas}
\end{equation}
Its linearization
\begin{equation}
 J_n:\ell^2(\mathcal P_n(B,M))\longrightarrow
          \ell^2(\mathcal E_{B^nM}),\qquad
 J_n|\pi\rangle=|\Phi_n(\pi)\rangle
 \label{lib05:jn}
\end{equation}
is unitary.
\end{theorem}

\begin{proof}
From a pair with sum \(B^nM\), the residue in
\eqref{lib05:inversa-etiquetas} recovers the last carry and a
unique preceding pair with sum \(B^{n-1}M\). Repeating this \(n\) times
produces a unique path whose input has sum \(M\). The forward
composition recovers the original endpoint, proving bijectivity and
the inverse. The count also follows from the inputs: the
\(M\) pairs with \(x_0\geq1\) each admit \(B^n\) words,
whereas \((0,M)\) admits only the all-zero
word. The total is \(MB^n+1\). The bijective transformation of
orthonormal bases proves the unitarity of \(J_n\).
\end{proof}

Once the depth is known, the readers on the final pair are
\begin{equation}
 c_j=\left\lfloor\frac{Y}{B^{n-j}}\right\rfloor\bmod B,
 \quad y_0=\left\lfloor\frac{Y}{B^n}\right\rfloor,
 \quad x_0=M-y_0,
 \qquad 1\leq j\leq n.
 \label{lib05:lectores-ruta}
\end{equation}
These functions are diagonal and recover the complete memory of
the finite operations. Compatibility with prolongation is
\begin{equation}
 \Phi_{n+1}(\pi,c)=L_c(\Phi_n(\pi)).
 \label{lib05:compatibilidad-prolongacion}
\end{equation}
Passing from \(n\) to \(n+1\) incorporates the new admissible
carry label. The spaces increase in dimension, and the
boundary label admits only \(c=0\); the domain of the
next step therefore retains admissibility instead of replacing it
with an unrestricted product with \(B\) digits.

\subsection{Fixed input and exact image}
\label{lib05:entrada-fija}

If an initial pair \((x_0,y_0)\) is fixed, with \(x_0\geq1\),
there are \(B^n\) paths, and their image is exactly
\begin{equation}
 \bigl\{(B^nx_0-C,B^ny_0+C):0\leq C\leq B^n-1\bigr\}.
 \label{lib05:imagen-fija}
\end{equation}
The result in \cref{lib03:cor-palabra} proves that the \(B^n\)
accumulated values are distinct and recover the respective words. The
linearization is unitary onto the subspace spanned by this
image. If \(x_0=0\), the image consists of a single label.

The cylinders of \cref{lib03:teo-cilindros} correspond to a
fixed input. The full domain of
\eqref{lib05:biyeccion-rutas} comprises all pairs with sum \(M\).
Both constructions have explicit inverses, but their images
and cardinalities differ. Arbitrary depth expresses the
validity of the proof for every \(n\); the limit of the readout and
its boundaries remain those established in
\cref{lib03:limite}.

\subsection{Bilateral coupling of refinement}
\label{lib05:bilateral}

Let \(R\) be a previously declared permutation of
\(\mathcal D_{B,M}\), linearized unitarily on
\(\mathscr H_D\). Define the two transports
\begin{equation}
 \mathcal B=J,\qquad \mathcal C=JR:
            \mathscr H_D\longrightarrow\mathscr H_E.
 \label{lib05:dos-transportes}
\end{equation}
The permutation is part of the specification of the procedure
and is fixed before an output readout is evaluated.
For \(a>0\), \(p=a^2+a+1\), \(N=(2a+1)p\),
\begin{equation}
 Q_-=J(aI-R),\qquad Q_+=J((a+1)I+R).
 \label{lib05:q-refinamiento}
\end{equation}

\begin{proposition}[Distribution and recovery of refined amplitudes]
\label{lib05:prop-bilateral}
The operators \eqref{lib05:q-refinamiento} satisfy
\begin{equation}
 (a+1)Q_-^*Q_-+aQ_+^*Q_+=NI_D,
 \label{lib05:balance}
\end{equation}
and the map
\begin{equation}
 W_a\psi=\frac1{\sqrt N}
    \bigl(\sqrt{a+1}\,Q_-\psi,\sqrt a\,Q_+\psi\bigr)
 \label{lib05:w}
\end{equation}
is isometric. For the unnormalized channels
\(u=Q_-\psi\), \(v=Q_+\psi\),
\begin{equation}
 J\psi=\frac{u+v}{2a+1},\qquad
 \psi=J^*\frac{u+v}{2a+1}.
 \label{lib05:inversa-bilateral}
\end{equation}
\end{proposition}

\begin{proof}
The two maps \(\mathcal B,\mathcal C\) are unitary.
The cross terms in the Gram operators of \(Q_-\) and \(Q_+\) have
coefficients \(-a\) and \(a+1\), respectively. The balance
weights cancel them, while the diagonal terms sum to
\(N\), as in \cref{lib04:teo-balance}. Normalization
proves the isometry. The sum of the channels is
\((2a+1)J\psi\), and the unitarity of \(J\) gives the inverse.
\end{proof}

The input is recovered directly from the normalized channels
by \(W_a^*\), with inverse norm one; this adjoint already incorporates
\(J^*\). The alternative formula
\eqref{lib05:inversa-bilateral} separates recovery of \(J\psi\)
by summation from application of \(J^*\), whose action on labels
is computed through residues. Reversible arithmetic refinement
and bilateral transport of amplitudes are thus composed.

\begin{proposition}[Compatibility in the refined space]
\label{lib05:prop-compatibilidad}
An unnormalized pair \((u,v)\) belongs to the image of
\(\psi\mapsto(Q_-\psi,Q_+\psi)\) if and only if
\begin{equation}
 \bigl[(a+1)I_E+JRJ^*\bigr]u
 =\bigl[aI_E-JRJ^*\bigr]v.
 \label{lib05:compatibilidad}
\end{equation}
\end{proposition}

\begin{proof}
The equality is equivalent to
\(JRJ^*(u+v)=av-(a+1)u\). If it holds, define
\(\psi=J^*(u+v)/(2a+1)\). Substituting into
\eqref{lib05:q-refinamiento}, this identity gives
\(Q_-\psi=u\) and \(Q_+\psi=v\). Necessity follows by
inserting the actual channels. All compositions act
on \(\mathscr H_E\), so their types are well defined.
\end{proof}

\subsection{Law of observables on the two channels}
\label{lib05:ley-lectores}

\begin{theorem}[Bilateral compression of a refined reader]
\label{lib05:teo-lector}
For \(F=F^*\) on \(\mathscr H_E\), write
\(F_J=J^*FJ\). Then
\begin{equation}
 W_a^*\operatorname{diag}(F,F)W_a
 =\frac{a(a+1)F_J+R^*F_JR}{p}.
 \label{lib05:lector-bilateral}
\end{equation}
The reader \(F_J\) is preserved exactly if and only if it commutes
with \(R\).
\end{theorem}

\begin{proof}
Expanding \((a+1)Q_-^*FQ_-+aQ_+^*FQ_+\) cancels the
mixed terms. The coefficient of \(\mathcal B^*F\mathcal B\)
is \(a(a+1)(2a+1)\); that of
\(\mathcal C^*F\mathcal C\) is \(2a+1\). Dividing by
\(N\), substituting \(\mathcal C=JR\), and using
\(F_J=J^*FJ\) proves the formula. Equality with \(F_J\)
is equivalent to \(R^*F_JR=F_J\), or to \([F_J,R]=0\) by
unitarity.
\end{proof}

In the case \(a=8\), the weights are \(72/73\) and \(1/73\).
The identity reader recovers the norm. A nonconstant reader
may change even when the norm is preserved; its variation is
computed by \eqref{lib05:lector-bilateral}.
The operator that preserves an original reader \(A\) exactly
on the full image is \(W_aAW_a^*\), whose image identity
is \(W_aW_a^*\). Reading \(\operatorname{diag}(F,F)\) corresponds
to the specific compression in the theorem and does not
automatically replace that transported observable.

\subsection{Arithmetic example in dimension one thousand and one}
\label{lib05:ejemplo}

Take \(B=10\), \(M=100\). The label spaces have
dimension \(1001\), and
\begin{equation}
 d_1=(73,27,1)\longmapsto e_1=(729,271).
 \label{lib05:ejemplo-etiqueta}
\end{equation}
The diagonal readers give
\begin{equation}
 f_E(e_1)=\frac{729}{1000}
 =\frac{73}{100}-\frac1{1000},\qquad
 g_E(e_1)=\frac{271}{1000},\qquad \gamma_E(e_1)=1.
 \label{lib05:ejemplo-lectores}
\end{equation}
Declare \(R\) to be the transposition of the admissible
labels \(d_1=(73,27,1)\) and \(d_2=(73,27,2)\), leaving
all others fixed. The second image is \(e_2=(728,272)\).
The transposition rule has been fixed on the labels before
any comparison with a constant.

For \(a=8\) and \(\psi=|d_1\rangle\),
\begin{equation}
 Q_-\psi=8|e_1\rangle-|e_2\rangle,\qquad
 Q_+\psi=9|e_1\rangle+|e_2\rangle.
 \label{lib05:ejemplo-canales}
\end{equation}
The squared norms are \(65\) and \(82\), and
\begin{equation}
 9\cdot65+8\cdot82=585+656=1241=17\cdot73.
 \label{lib05:ejemplo-balance}
\end{equation}
The sum of the channels is \(17|e_1\rangle\). The inverse
\eqref{lib05:inversa-bilateral} and the residue recover
\(|d_1\rangle\), including the carry one.

If \(F=M_{f_E}\) is evaluated diagonally on both normalized
channels, \cref{lib05:teo-lector} gives
\begin{equation}
 \begin{split}
 \langle W_8\psi,\operatorname{diag}(F,F)W_8\psi\rangle
 &=\frac{72\cdot729+728}{73\cdot1000}\\
 &=\frac{729}{1000}-\frac1{73000}.
 \end{split}
 \label{lib05:ejemplo-fraccion}
\end{equation}
The complementary readout is one minus this quantity. For the
diagonal carry reader, the result is
\begin{equation}
 \frac{72\cdot1+2}{73}=\frac{74}{73}.
 \label{lib05:ejemplo-acarreo}
\end{equation}
The original carry remains exactly recoverable: the
observable \(W_8M_cW_8^*\) has eigenvalue one on
\(W_8|d_1\rangle\). The value \(74/73\) is the response of the
diagonal reader applied to both channels, whose relative transport
is specified by \eqref{lib05:lector-bilateral}. The example
distinguishes preservation of information from variation in a particular
readout.

\subsection{Relation to the generated register and domains of application}
\label{lib05:dominios}

The unitary \(J\) acts on \(BM+1\) labels, whereas
the realization of \(K\) has twelve coordinates. In the example,
\(1001\) counts refined pairs with sum one thousand, and twelve counts
windows of the produced register. Both dimensions retain their
specific combinatorial origin.

A particular composition between these two carriers can be formulated
through a materially specified selector of twelve distinct paths, with its
provenance and operators. If this selector provides an
isometry
\(J_{12}:\mathbb R^{12}\to\ell^2(\mathcal D_{B,M})\), then
\(JJ_{12}\) is isometric, and the intertwining of readers is
checked on it. This is the hypothesis of that additional
composition, not an implicit identification based on the number of
coordinates. The bilateral maps already defined directly
on \(K\) in \cref{lib04:bilateral} retain their full domains
without using that additional selector.

The result of this section brings together an arithmetic bijection, its
unitary lift, intertwined diagonal readers, recovery of paths at
every finite depth, and an explicit bilateral composition.
The path domain is the one defined in
\eqref{lib05:biyeccion-rutas}. Applications to histories
filtered by TPK retain their survival and prolongation
rules; the label correspondence proves the declared
transport here without replacing those operators with an unrestricted
arithmetic domain.
